\mxunit{expertes-c06}{Divisibilité et congruences}
\mxsubtitle{Divisibilité dans $\Z$, division euclidienne, congruences}

\begin{objectives}
Utiliser la divisibilité dans $\Z$ et ses propriétés ; effectuer une division euclidienne ; calculer avec des congruences ; établir un critère de divisibilité ; étudier des restes de puissances.
\end{objectives}
\begin{prerequisites}
Multiples, diviseurs, nombres premiers (seconde) ; raisonnement par récurrence ; partie entière d'un réel.
\end{prerequisites}

\section{Automatismes}
\begin{fluency}[Sans calculatrice, cinq minutes]
\begin{enumerate}
\item Donner la liste des diviseurs positifs de $36$.
\item Le nombre $n(n+1)$ est-il toujours pair ? Justifier.
\item Quel est le chiffre des unités de $7^4$ ?
\item Simplifier la fraction $\frac{84}{126}$.
\item Le nombre $\num{2026}$ est-il divisible par $3$ ? par $2$ ?
\item Calculer $17\times23$.
\end{enumerate}
\end{fluency}
\begin{solution}[Automatismes]
\begin{enumerate}
\item $1$, $2$, $3$, $4$, $6$, $9$, $12$, $18$, $36$.
\item Oui : l'un des deux entiers consécutifs est pair.
\item $7^2=49$, $7^4=\num{2401}$ : le chiffre des unités est $1$.
\item $\frac{84}{126}=\frac{2\times42}{3\times42}=\frac23$.
\item La somme des chiffres vaut $10$, non multiple de $3$ : non. Il est pair : oui.
\item $391$.
\end{enumerate}
\end{solution}

\mxnewpage
\section{Activité d'introduction}
\begin{activity}[Quel jour sera-t-il ? \mxchip{35 min}]
Le 17 octobre 2026 est un samedi. On numérote les jours de la semaine de $0$ (samedi) à $6$ (vendredi) : dimanche porte le numéro $1$, lundi le numéro $2$, etc.
\begin{enumerate}
\item Quel jour de la semaine sera-t-il $7$ jours plus tard ? $14$ jours ? $700$ jours ?
\item Quel jour sera-t-il $100$ jours plus tard ? Expliquer comment obtenir la réponse sans compter les jours un à un.
\item Même question pour $\num{1000}$ jours.
\item L'année 2027 n'est pas bissextile. Quel jour de la semaine sera le 17 octobre 2027 ?
\item Le 1\textsuperscript{er} janvier 2027 arrive $76$ jours après le 17 octobre 2026. Quel jour sera-t-il ?
\item Deux durées $a$ et $b$, en jours, mènent au même jour de la semaine. Que peut-on dire de $a-b$ ? Formuler une condition générale.
\end{enumerate}
\par\emph{Observer n'est pas démontrer : les propriétés conjecturées ici sont établies dans le cours.}
\end{activity}
\begin{solution}[Activité]
\begin{enumerate}
\item Samedi dans les trois cas, car $7$, $14$ et $700$ sont des multiples de $7$.
\item $100=7\times14+2$ : on avance de $2$ jours, c'est un lundi. Seul compte le reste de la division par $7$.
\item $\num{1000}=7\times142+6$ : vendredi.
\item $365=7\times52+1$ : dimanche.
\item $76=7\times10+6$ : vendredi.
\item $a$ et $b$ ont le même reste dans la division par $7$, c'est-à-dire que $a-b$ est un multiple de $7$. On dira que $a$ et $b$ sont congrus modulo $7$.
\end{enumerate}
\end{solution}
\begin{teachernote}[Conduite de l'activité]
Faire vérifier les réponses à l'aide d'un calendrier ou d'un tableur pour la question 4, en signalant que la vérification sur un exemple ne remplace pas le raisonnement. Prolongement : expliquer pourquoi le 17 octobre 2030 est un jeudi ($\num{1461}$ jours, dont un 29 février, et $\num{1461}=7\times208+5$).
\end{teachernote}

\mxnewpage
\section{Notes de cours}
\subsection{Divisibilité dans $\Z$}
\begin{definition}[Diviseur, multiple]
Soit $a$ et $b$ deux entiers relatifs. On dit que $a$ \emph{divise} $b$, et on note $a\mid b$, s'il existe un entier relatif $k$ tel que $b=ka$. On dit aussi que $a$ est un \emph{diviseur} de $b$, ou que $b$ est un \emph{multiple} de $a$.
\end{definition}
\begin{theorem}[Propriétés de la divisibilité]
Soit $a$, $b$, $c$ des entiers relatifs.
\begin{enumerate}
\item Si $a\mid b$ et $b\mid c$, alors $a\mid c$.
\item Si $a\mid b$ et $a\mid c$, alors $a\mid bu+cv$ pour tous entiers relatifs $u$ et $v$.
\item Si $a\mid b$ et $b\neq0$, alors $\abs a\le\abs b$.
\end{enumerate}
\end{theorem}
\begin{proof}
1. $b=ka$ et $c=k'b$ donnent $c=(k'k)a$. 2. $b=ka$ et $c=k'a$ donnent $bu+cv=(ku+k'v)a$. 3. $b=ka$ avec $k\neq0$, donc $\abs k\ge1$ et $\abs b=\abs k\abs a\ge\abs a$.
\end{proof}
\begin{theorem}[Division euclidienne]
Soit $a\in\Z$ et $b\in\N$, $b\ge1$. Il existe un unique couple d'entiers $\pt{q;r}$ tel que $a=bq+r$ et $0\le r<b$. On appelle $q$ le \emph{quotient} et $r$ le \emph{reste}.
\end{theorem}
\begin{proof}
\emph{Existence.} Posons $q=\left\lfloor\frac ab\right\rfloor$, le plus grand entier inférieur ou égal à $\frac ab$. Alors $q\le\frac ab<q+1$, donc $bq\le a<bq+b$, et $r=a-bq$ vérifie $0\le r<b$.\par
\emph{Unicité.} Si $a=bq+r=bq'+r'$ avec $0\le r,r'<b$, alors $b(q-q')=r'-r$ et $-b<r'-r<b$. Le seul multiple de $b$ strictement compris entre $-b$ et $b$ est $0$, donc $r=r'$ puis $q=q'$.
\end{proof}
\begin{pitfall}
Le reste est toujours positif : $-47=5\times(-10)+3$, donc la division de $-47$ par $5$ a pour quotient $-10$ et pour reste $3$, et non $-9$ et $-2$.
\end{pitfall}

\subsection{Congruences}
\begin{definition}[Congruence modulo $n$]
Soit $n\ge2$ un entier. Deux entiers relatifs $a$ et $b$ sont \emph{congrus modulo $n$}, ce que l'on note $a\equiv b\pmod n$, si $n\mid a-b$.
\end{definition}
\begin{theorem}[Congruence et reste]
$a\equiv b\pmod n$ si et seulement si $a$ et $b$ ont le même reste dans la division euclidienne par $n$. En particulier, tout entier est congru modulo $n$ à son reste, qui appartient à $\ent{0}{n-1}$.
\end{theorem}
\begin{proof}
Écrivons $a=nq+r$ et $b=nq'+r'$ avec $0\le r,r'<n$. Alors $a-b=n(q-q')+(r-r')$ et $-n<r-r'<n$. Donc $n\mid a-b$ si et seulement si $n\mid r-r'$, si et seulement si $r=r'$.
\end{proof}
\begin{theorem}[Compatibilité avec les opérations]
Si $a\equiv b\pmod n$ et $c\equiv d\pmod n$, alors
\[a+c\equiv b+d,\qquad a-c\equiv b-d,\qquad ac\equiv bd,\qquad a^k\equiv b^k\pmod n\]
pour tout entier naturel $k$. La relation est aussi transitive : si $a\equiv b$ et $b\equiv c\pmod n$, alors $a\equiv c\pmod n$.
\end{theorem}
\begin{proof}
Somme : $(a+c)-(b+d)=(a-b)+(c-d)$ est somme de deux multiples de $n$. Produit : $ac-bd=a(c-d)+d(a-b)$, multiple de $n$. Puissance : récurrence sur $k$ à partir du produit. Transitivité : $a-c=(a-b)+(b-c)$.
\end{proof}
\begin{pitfall}
On ne simplifie pas une congruence comme une égalité : $2\times3\equiv2\times1\pmod4$, mais $3\not\equiv1\pmod4$. La division n'est possible que dans des cas précis (chapitre 7).
\end{pitfall}
\begin{method}[Reste d'une grande puissance]
On cherche une puissance congrue à $1$ (ou à $-1$), puis on utilise la division euclidienne de l'exposant. Exemple : $2^3=8\equiv1\pmod7$ et $\num{2026}=3\times675+1$, donc $2^{2026}=\left(2^3\right)^{675}\times2\equiv1^{675}\times2=2\pmod7$.
\end{method}
\begin{method}[Tableau de congruences]
Pour étudier une expression modulo $n$, on dresse un tableau donnant sa valeur modulo $n$ pour chaque reste possible de $x$. Exemple : modulo $4$, $x\equiv0,1,2,3$ donne $x^2\equiv0,1,0,1$. Un carré est donc congru à $0$ ou $1$ modulo $4$.
\end{method}
\begin{method}[Critère de divisibilité par $9$]
Si $N=\sum_{k=0}^{p}a_k10^k$ s'écrit avec les chiffres $a_0,\dots,a_p$, comme $10\equiv1\pmod9$, on a $10^k\equiv1$ et $N\equiv a_0+a_1+\dots+a_p\pmod9$. Un entier est divisible par $9$ si et seulement si la somme de ses chiffres l'est.
\end{method}

\mxnewpage
\section{Exercices}
\begin{exercise}[Divisions euclidiennes]
\begin{enumerate}
\item Effectuer la division euclidienne de $\num{2026}$ par $17$.
\item Effectuer la division euclidienne de $-47$ par $5$, puis de $-35$ par $4$.
\item Déterminer les entiers naturels $n$ dont le quotient dans la division par $6$ est égal au reste.
\end{enumerate}
\end{exercise}
\begin{solution}
\begin{enumerate}
\item $\num{2026}=17\times119+3$ : quotient $119$, reste $3$.
\item $-47=5\times(-10)+3$ et $-35=4\times(-9)+1$.
\item $n=6q+q=7q$ avec $0\le q<6$ : $n\in\ens{0;7;14;21;28;35}$.
\end{enumerate}
\end{solution}

\begin{exercise}[Restes de puissances]
\begin{enumerate}
\item Déterminer le reste de $3^{100}$ dans la division par $5$.
\item Déterminer le reste de $5^{2026}+2^{2026}$ dans la division par $3$.
\end{enumerate}
\end{exercise}
\begin{solution}
\begin{enumerate}
\item $3^4=81\equiv1\pmod5$, donc $3^{100}=\left(3^4\right)^{25}\equiv1\pmod5$ : le reste est $1$.
\item $5\equiv2\equiv-1\pmod3$, donc $5^{2026}+2^{2026}\equiv2\times(-1)^{2026}=2\pmod3$ : le reste est $2$.
\end{enumerate}
\end{solution}

\begin{exercise}[Sommes de deux carrés]
\begin{enumerate}
\item Démontrer que pour tout entier $x$, $x^2\equiv0$ ou $x^2\equiv1\pmod4$.
\item En déduire qu'un entier congru à $3$ modulo $4$ n'est pas la somme de deux carrés d'entiers.
\item L'entier $\num{2027}$ peut-il s'écrire $x^2+y^2$ avec $x$, $y$ entiers ?
\end{enumerate}
\end{exercise}
\begin{solution}
\begin{enumerate}
\item Tableau : $x\equiv0,1,2,3\pmod4$ donne $x^2\equiv0,1,4,9$, c'est-à-dire $0,1,0,1\pmod4$.
\item $x^2+y^2$ est congru à $0+0$, $0+1$ ou $1+1$, donc à $0$, $1$ ou $2$ modulo $4$, jamais à $3$.
\item $\num{2027}=4\times506+3$ : non.
\end{enumerate}
\end{solution}

\begin{exercise}[Critère de divisibilité par $11$]
\begin{enumerate}
\item Démontrer que $10\equiv-1\pmod{11}$, puis que $10^k\equiv(-1)^k\pmod{11}$.
\item En déduire qu'un entier est divisible par $11$ si et seulement si la somme alternée de ses chiffres, en partant des unités, $a_0-a_1+a_2-\dots$, l'est.
\item L'entier $\num{918082}$ est-il divisible par $11$ ?
\end{enumerate}
\end{exercise}
\begin{solution}
\begin{enumerate}
\item $10-(-1)=11$. La compatibilité avec les puissances donne $10^k\equiv(-1)^k$.
\item $N=\sum a_k10^k\equiv\sum a_k(-1)^k\pmod{11}$ ; $N$ et la somme alternée ont le même reste.
\item $2-8+0-8+1-9=-22$, multiple de $11$ : oui ($\num{918082}=11\times\num{83462}$).
\end{enumerate}
\end{solution}

\begin{exercise}[Divisibilités pour tout entier]
Démontrer que pour tout entier naturel $n$ :
\begin{enumerate}
\item $4^n-1$ est divisible par $3$ ;
\item $n^3-n$ est divisible par $6$.
\end{enumerate}
\end{exercise}
\begin{solution}
\begin{enumerate}
\item $4\equiv1\pmod3$, donc $4^n\equiv1$ et $3\mid4^n-1$.
\item Tableau modulo $6$ : pour $n\equiv0,1,2,3,4,5$, $n^3-n\equiv0,0,6,24,60,120$, tous multiples de $6$. Donc $n^3-n\equiv0\pmod6$ pour tout $n$. Remarque : $n^3-n=(n-1)n(n+1)$ est un produit de trois entiers consécutifs ; on verra au chapitre 7 pourquoi être divisible par $2$ et par $3$ suffit.
\end{enumerate}
\end{solution}

\begin{exercise}[Équations de congruence]
\begin{enumerate}
\item Résoudre $3x\equiv5\pmod7$. On cherchera d'abord un entier $u$ tel que $3u\equiv1\pmod7$.
\item Résoudre $x^2\equiv2\pmod7$ à l'aide d'un tableau.
\end{enumerate}
\end{exercise}
\begin{solution}
\begin{enumerate}
\item $3\times5=15\equiv1\pmod7$. En multipliant par $5$ : $3x\equiv5$ implique $x\equiv25\equiv4\pmod7$ ; réciproquement $3\times4=12\equiv5$. Les solutions sont les entiers $x=4+7k$, $k\in\Z$.
\item Pour $x\equiv0,1,\dots,6\pmod7$ : $x^2\equiv0,1,4,2,2,4,1$. Les solutions sont les entiers congrus à $3$ ou $4$ modulo $7$.
\end{enumerate}
\end{solution}

\begin{exercise}[Les deux derniers chiffres]
Déterminer les deux derniers chiffres de l'écriture décimale de $7^{2026}$.
\end{exercise}
\begin{solution}
On travaille modulo $100$. $7^2=49$ et $7^4=\num{2401}\equiv1\pmod{100}$. Comme $\num{2026}=4\times506+2$, $7^{2026}\equiv7^2=49\pmod{100}$ : l'écriture se termine par $49$.
\end{solution}

\mxnewpage
\section{Problème}
\begin{problem}[La clé de contrôle d'un code ISBN-13]
Un code ISBN-13 est une suite de treize chiffres $a_1a_2\dots a_{13}$. On pose
\[S=a_1+3a_2+a_3+3a_4+\dots+3a_{12}+a_{13},\]
les poids $1$ et $3$ alternant. Le code est valide si et seulement si $S\equiv0\pmod{10}$ ; le dernier chiffre $a_{13}$, appelé clé, est choisi pour cela.
\begin{enumerate}
\item On considère le numéro fictif 978-2-0601-2345-$a_{13}$. Calculer la clé $a_{13}$.
\item On suppose qu'un seul chiffre $a$ est remplacé par un chiffre $a'\neq a$.
\begin{enumerate}
\item Démontrer que $S$ est modifié de $w(a'-a)$, avec $w\in\ens{1;3}$.
\item Vérifier que $7\times3\equiv1\pmod{10}$ et en déduire que $3(a'-a)\equiv0\pmod{10}$ implique $a'-a\equiv0\pmod{10}$.
\item Conclure : une erreur sur un seul chiffre est toujours détectée.
\end{enumerate}
\item On suppose que deux chiffres voisins distincts $a$ et $b$ sont échangés.
\begin{enumerate}
\item Démontrer que $S$ est modifié de $2(b-a)$ ou de $2(a-b)$.
\item Déterminer à quelle condition sur $a$ et $b$ l'erreur n'est pas détectée. Donner un exemple.
\end{enumerate}
\end{enumerate}
\end{problem}
\begin{solution}[Problème]
\begin{enumerate}
\item Les douze premiers chiffres $9,7,8,2,0,6,0,1,2,3,4,5$ donnent $9+21+8+6+0+18+0+3+2+9+4+15=95$. Il faut $95+a_{13}\equiv0\pmod{10}$ avec $0\le a_{13}\le9$ : $a_{13}=5$.
\item
\begin{enumerate}
\item Le terme $wa$ devient $wa'$, les autres termes ne changent pas.
\item $21\equiv1\pmod{10}$. Si $3(a'-a)\equiv0$, en multipliant par $7$ : $a'-a\equiv21(a'-a)\equiv0\pmod{10}$. Le cas $w=1$ est immédiat.
\item $a'-a$ est un multiple de $10$ compris entre $-9$ et $9$, donc nul : contradiction avec $a'\neq a$. La variation $w(a'-a)$ de la somme n'est pas multiple de $10$ ; le code d'origine étant valide (somme multiple de $10$), la nouvelle somme n'est plus multiple de $10$ et l'erreur est détectée.
\end{enumerate}
\item
\begin{enumerate}
\item Les deux chiffres ont des poids $1$ et $3$. La contribution $a+3b$ devient $b+3a$ (ou l'inverse), soit une variation de $2(a-b)$ ou de $2(b-a)$.
\item L'erreur n'est pas détectée si et seulement si $10\mid2(a-b)$, c'est-à-dire $5\mid a-b$ : $\abs{a-b}=5$. Exemple : échanger les chiffres voisins $1$ et $6$ ne modifie pas $S$ modulo $10$.
\end{enumerate}
\end{enumerate}
\end{solution}
\begin{teachernote}[Statut du problème]
Les codes de contrôle (ISBN, RIB, codes-barres) sont un « problème possible » du programme. La question 2b utilise un inverse modulo $10$ : la notion sera formalisée au chapitre 7. Le numéro de la question 1 est donné à titre d'exemple, pour le calcul.
\end{teachernote}
